\subsection{Further indivisibility inputs and the selected quadratic value}

\paragraph{KT-INDIVISIBILITY-LITERATURE-SCOPE.}
The missing modular-symbol assertion concerns the specified
$j=2$ value of the integral Eisenstein symbol, hence the same
negative fundamental discriminant in
$\beta_3(\chi_D)\beta_b(\chi_D)$.
It is a noncentral value: $L(E,\chi_D,3)$ with
$k=(p+9)/2\geq20$.  The following source checks record why
the particular imported results inspected do not prove that
this value is a unit.  They do not claim that no applicable
theorem can exist.

Heumann--Vatsal, \emph{Modular symbols, Eisenstein series,
and congruences}, J.~Th\'eor.~Nombres Bordeaux \textbf{26} (2014),
709--756, Theorem~5.1, pp.~747--748, gives nonvanishing with
characters of growing odd prime-power order and uses $j=0$.
The already checked Corollary~4.7 gives our quadratic
coefficient formula but no nonzero selected quadratic value.
Theorem~5.2 transfers congruences without adding that unit.
The full Theorem~5.1 proof was read in the separate detector
audit.

Bruinier, \emph{Nonvanishing modulo $\ell$ of Fourier
coefficients of half-integral weight modular forms},
Duke Math.~J.~\textbf{98} (1999), 595--611, Theorem~1,
requires the residue prime to be prime to the auxiliary
Legendre prime times one less than that prime.  It therefore
does not impose that condition at our chosen $q\equiv1\pmod p$.
Another auxiliary prime is not excluded.  More basically,
Theorems~2--4 concern one half-integral weight eigenform;
Theorem~6 and its ensuing odd-index remark concern one
Dirichlet value.  No eigenform with fundamental coefficients
equal to our product has been constructed here.  The exact
statements on author pp.~6 and 10--13 were read, together
with the BJKOSV survey, pp.~6--7 and 13--14; no foundational
whole-proof audit is claimed.

Ahlgren--Boylan, \emph{Central critical values of modular
$L$-functions and coefficients of half-integral weight modular
forms modulo $\ell$}, Amer.~J.~Math.~\textbf{129} (2007),
429--454, author version dated 29 June 2005,
Theorem~1.1, p.~2,
\url{https://people.math.sc.edu/boylan/reprints/THETAFINALNEW.pdf},
concerns central quadratic values of a cuspidal newform of
even weight and odd squarefree level.  Its hypothesis that
a minimum valuation is attained for two discriminants does
not say that this minimum is zero in our normalization.
Only the complete theorem statement and hypotheses were
retrieved; the later PDF download failed, so no local copy
or whole-proof reading is claimed.

Kayath--Lane--Neifeld--Ni--Xue, \emph{Subspaces spanned by
eigenforms with nonvanishing twisted central $L$-values},
Canadian J.~Math.,
\url{https://doi.org/10.4153/S0008414X25101697},
Theorem~1.2 and Proposition~2.4, pp.~7--8 of the author
version, give a complex spanning construction that divides
by selected nonzero complex central values.  The full
Proposition~2.4 proof was read.  Such division does not
establish saturation of an integral lattice at $p$;
the chosen values could be divisible by $p$.  Vatsal,
\emph{Canonical periods and congruence formulae},
Theorem~0.3, author p.~3, and its Eisenstein discussion
on p.~14, applies in its stated quadratic application to
ordinary elliptic curves with rational three-torsion and
squarefree conductor.  These inspected statements do not
give the arbitrary-prime, higher-weight, specified
noncentral coefficient required here.

\paragraph{KT-KIMURA-DISPLAYED-PROOF-SCOPE.}
I.~Kimura, \emph{Indivisibility of special values of zeta
functions associated to real quadratic fields},
RIMS K\^oky\^uroku Bessatsu B4 (2007), Theorem~3.1
and its full proof, printed pp.~11--15,
\url{https://www.kurims.kyoto-u.ac.jp/~kenkyubu/bessatsu/open/B4/pdf/B04_001.pdf},
was checked for a possible $p$-ramified fixed-weight input.
Two displayed steps need additional justification for
that application.  This is a gap in the displayed
argument as used here, not a counterexample to its theorem.

The form in the proof is first depleted at $p$.
Its surviving Fourier indices are prime to $p$.
The later sieves and $U_\ell,V_\ell$ operations with
$\ell\ne p$ preserve this property: multiplying or
dividing an index by $\ell$ cannot change its divisibility
by $p$.  Thus the final nonzero coefficient $H(r,n\ell)$
has $p\nmid n\ell$.  Writing $n\ell=Df^2$ then gives
$p\nmid D$.  This coefficient construction does not as
written yield the stated $p$-ramified discriminant
$D_\ell=d_\ell\ell p$.

There is also a separate coefficient-comparison issue.
Write $a=\ell^{r-1}$, $b=\ell^{2r-1}$ and $c=1-a+b$.
Equation~(3.6) is $U_\ell G_r=cV_\ell G_r$.
At the chosen index $D_0\ell^3$, the normalized
$U$ coefficient is $1+b+b^2-a-ab$ and the
$V$ coefficient is $c$.  Hence the comparison required is
\[
 (1+b+b^2-a-ab)-c^2
   =a(1-a+b)-b
   =a(1-\ell^r)(1-a),
 \tag{KS.1}
\]
where the last equality uses $b=a\ell^r$.
The printed subtraction instead compares the $U$
coefficient with $c$; it addresses $U\ne V$, not
$U\ne cV$.  Adding $\ell^{r-1}\ne1\pmod p$
repairs this coefficient comparison when an admissible
residue can be chosen.  It does not repair the distinct
support issue.  No result in this paper relies on the
unjustified $p$-ramified conclusion.

K.~Ono, \emph{Indivisibility of class numbers of real
quadratic fields}, Compositio Math.~\textbf{119} (1999),
1--11, Theorem~2, and D.~Byeon, \emph{Indivisibility of
special values of Dedekind zeta functions of real quadratic
fields}, Acta Arith.~\textbf{109.3} (2003),
Theorems~1--2, were also inspected for individual
quadratic-field indivisibility.  Only their indicated
statements and initial construction were read; neither
was subjected to a whole-proof audit in this attempt.
Their precise retrieval records, and the fuller bounded
quadratic-symbol search with its source-by-source scope,
are retained in
\texttt{main\_gate/indivisibility\_source\_notes.md} and
\texttt{lean\_chain/QUADRATIC\_SYMBOL\_NONVANISHING\_SCOPE.md}.

These checks leave simultaneous auxiliary-unit existence,
the nonessentially-selfdual source exclusion, and the
original strict rational class component open.
